\section{引言：从 2D 先验到跨模态生成}

本节课是 KAIST CS492D（Diffusion Models and Their Applications）的最后一讲。Minhyuk Sung 教授在前几讲中已经讨论了扩散模型的基础理论（DDPM、DDIM、连续时间扩散、SDE/ODE 转换、Flow Matching）以及推理时引导（inference-time guidance）等技术。本节课聚焦一个核心问题：\textbf{如何利用预训练的 2D 图像扩散模型生成训练数据中从未出现过的其他模态数据}——尤其是 3D 内容。

\begin{knowledgebox}{课程回顾：从引导到蒸馏}
前几讲中，我们讨论了如何在推理阶段注入引导信号（guidance），使扩散模型生成满足特定约束的样本。本节课的 Score Distillation Sampling（SDS）则更进一步：它不是在已有数据域中做引导生成，而是将预训练扩散模型的知识``蒸馏''到一个全新的数据域中。
\end{knowledgebox}

\subsection{动机：数据规模的鸿沟}

扩散模型在图像和视频生成上取得了巨大成功，核心原因之一是训练数据的规模——图像数据集已达数十亿级别（如 LAION-5B 约 50 亿张图像）。然而，对于其他数据类型，数据规模远远不及：

\begin{table}[H]
\centering
\begin{tabular}{lcc}
\toprule
\textbf{数据类型} & \textbf{最大公开数据集规模} & \textbf{类比} \\
\midrule
2D 图像 & $\sim$50 亿 & 世界人口 \\
3D 模型 & $\sim$1000 万 & 首尔人口 \\
分子结构 & 数百万级 & 更少 \\
运动/动作数据 & 数十万级 & 很少 \\
\bottomrule
\end{tabular}
\caption{不同数据类型的公开数据集规模对比}
\end{table}

\begin{importantbox}{核心问题}
图像扩散模型因数十亿级训练数据而具备了强大的生成能力与多样性。但对于 3D 模型、矢量图、CAD、分子结构等领域，我们很难获得同等规模的训练数据。\textbf{能否借用图像扩散模型已学到的知识来生成这些跨模态数据？}
\end{importantbox}

这个问题的实际意义非常广泛：医学影像（用少量样本训练诊断模型）、工业设计（CAD 模型生成）、电商（3D 产品展示）等场景都面临数据不足的困境。

\subsection{本章小结}

本节引出了 Score Distillation 的核心动机：预训练图像扩散模型拥有在数十亿图像上学到的丰富先验知识，SDS 提供了一种方法将这些知识迁移到 3D 等其他领域，即便目标域没有大规模训练数据。

\newpage
